% !TeX program = pdflatex
% corpus/docs/ciencia_dragao.tex — ENSINAR CIÊNCIA pela cadeia do bit ao dragão.
%
% A solução dos problemas mapeados NÃO mora na cabeça da assistente nem em
% strings do conversa.c: mora AQUI, em ticks (passo a passo). A base só guarda
% @TEX este ficheiro; responde() lê o corpo AO VIVO e devolve em LaTeX.
%
% Cadeia (papers):
%   F_2 → F_8≡W_8 → ℕ (naturais) → ℤ (inteiros) → ℚ (racionais) → ℝ (reais)
%   → realização Heighway / aranha / inversa (arquitetura §aranha).
%
% Membrana: entrada ASCII (desdobra); saída LaTeX (veste). Ver §membrana.
%
\documentclass[11pt,a4paper]{article}
\usepackage[utf8]{inputenc}
\usepackage[T1]{fontenc}
\usepackage[portuguese]{babel}
\usepackage{amsmath,amssymb,amsthm}
\usepackage[margin=2.4cm]{geometry}
\usepackage{booktabs}
\usepackage{xcolor}
\usepackage[hidelinks]{hyperref}

\newtheorem{definicao}{Definição}
\newtheorem{proposicao}{Proposição}
\newtheorem{observacao}{Observação}

\newcommand{\code}[1]{\texttt{\small #1}}
\newcommand{\Word}{\mathcal{W}}
\newcommand{\medido}[1]{\par\medskip\noindent{\small\color{gray}\textbf{Medido.} #1}\par\medskip}

\begin{document}
\title{Ciência pelo dragão\\
\large cadeia do bit $\to$ Heighway --- ticks de resolução}
\author{Tiffany --- corpus vivo}
\date{}
\maketitle

\begin{abstract}
\noindent Ensinar ciência aqui é \textbf{seguir a cadeia} e resolver em
\textbf{ticks}: cada tick nomeia o que fez. A geometria final é a curva do
dragão (Heighway): $G>1$ na base, levantamento $\tilde\pi$ com
$\tilde G\equiv1$. Fontes: \code{papers/naturais.tex},
\code{inteiros.tex}, \code{racionais.tex}, \code{reais.tex},
\code{arquitetura.tex} §aranha. Nada se inventa fora destes papers --- este
documento \emph{ordena} os passos e firma a membrana LaTeX$\leftrightarrow$ASCII.
\end{abstract}

%----------------------------------------------------------------------
\section{a cadeia do basico ao dragao}\label{sec:cadeia}
\[
\boxed{\;
\mathrm{GF}(2)
\;\xrightarrow{\ \code{naturais}\ }\;
\Word_8
\;\xrightarrow{\ \mathrm{transporte}\ }\;
\mathbb{N}
\;\xrightarrow{\ \code{inteiros}\ }\;
\mathbb{Z}
\;\xrightarrow{\ \code{racionais}\ }\;
\mathbb{Q}
\;\xrightarrow{\ \code{reais}\ }\;
\mathbb{R}
\;\xrightarrow{\ \pi\ \mathrm{Heighway}\ }\;
G
\;\xrightarrow{\ \tilde\pi\ }\;
\tilde G\equiv1
\;}
\]
\noindent\textbf{Tick de leitura.} Da esquerda para a direita \emph{constrói-se};
da direita para a esquerda \emph{projecta-se}. O dragão não inventa a
aritmética --- realiza a órbita discreta $I$ que a álgebra já emitiu.

\section{como se ensina --- os ticks}
Cada resolução tem seis slots (o relógio da casa):
\begin{enumerate}\itemsep2pt
\item \textbf{DEFINIÇÃO} --- o objecto e o corpo onde vive;
\item \textbf{TRADUÇÃO} --- LaTeX $\leftrightarrow$ ASCII (membrana);
\item \textbf{OPERAÇÃO} --- a lei que se aplica;
\item \textbf{DEMONSTRAÇÃO} --- o passo com resíduo~$0$;
\item \textbf{CONTRAEXEMPLO} --- se falhar, onde falha;
\item \textbf{VOLTA} --- dual / projecção / verificação.
\end{enumerate}
Operacionalmente: \code{conversa <base> tick "\ldots"} ou a cascata
\code{resolve\_*} em \code{banco/conversa.c} (cada \texttt{tique7} é um slot
ao vivo). A prosa canónica das resoluções migra para \code{corpus/docs/}:
conservação métrica CM1--CM8 em \code{conservacao\_metrica.tex}.

%----------------------------------------------------------------------
\section{membrana latex ascii}\label{sec:membrana}
A interface padrão é LaTeX; o motor resolve em ASCII desdobrado; a resposta
\textbf{veste-se} de novo.

\begin{center}\small
\begin{tabular}{@{}lll@{}}
\toprule
sentido & operação & exemplo \\
\midrule
entrada ($+1$) & \code{latex\_desdobra} / \code{MEMBRANA}$(+1)$
  & \verb|\frac{1}{2}| $\to$ \texttt{(1)/(2)} \\
saída ($-1$) & \code{veste\_valor} / \code{MEMBRANA}$(-1)$
  & $v/q$ $\to$ \verb|\frac{v}{q}| \\
ASCII pontos & \code{lib/ascii.h} & $128=|\mathbb{P}^{1}(\mathbb{F}_{127})|$ \\
\bottomrule
\end{tabular}
\end{center}

\noindent\textbf{Regra.} Pergunta com \verb|\frac|, \verb|$| ou \verb|^{}| ---
desdobra na entrada. Resposta numérica da fita --- veste em LaTeX. Corpo lido
de \texttt{@TEX} --- \textbf{preserva} o LaTeX do ficheiro (não se limpa a matemática).

\subsection*{tabela de conversao}
\begin{center}\small
\begin{tabular}{@{}ll@{}}
\toprule
LaTeX & ASCII (fita) \\
\midrule
\verb|\frac{a}{b}| & \texttt{(a)/(b)} \\
\verb|\sqrt{x}| / \verb|\sqrt{x}| & \texttt{raiz (x)} \\
\verb|x^{2}| & \texttt{x\^{}2} \\
\verb|\cdot| / \verb|\times| & \texttt{x} \\
\verb|\mathbb{N},\mathbb{Z},\mathbb{Q},\mathbb{R}| & \texttt{N,Z,Q,R} (nome) \\
\verb|\oplus| & \texttt{xor} / soma em $\mathrm{GF}(2)$ \\
\bottomrule
\end{tabular}
\end{center}

%----------------------------------------------------------------------
\section{tick bit --- GF(2)}\label{sec:tick-bit}
\textbf{DEFINIÇÃO.} $F_1=\mathrm{GF}(2)=\{0,1\}$ com $1+1=0$.

\textbf{TRADUÇÃO.} \verb|1+1| em ASCII é a conta; em LaTeX
$1+1$ no corpo $\mathrm{GF}(2)$.

\textbf{OPERAÇÃO.} Em $\mathrm{GF}(2)$ a soma é XOR: $a\oplus b$.

\textbf{DEMONSTRAÇÃO.} $1\oplus1=0$ --- resíduo~$0$ na tabuada.

\textbf{CONTRAEXEMPLO.} Em $\mathbb{Z}$, $1+1=2\neq0$ --- o corpo tem de estar dito.

\textbf{VOLTA.} Dual em $\mathrm{GF}(2)$: $-x=x$.

\section{tick naturais --- transporte}\label{sec:tick-N}
\textbf{DEFINIÇÃO.} $\Word_8=\{0,\ldots,255\}$; $\mathbb{N}$ pela torre
$F_{2w}=F_w+F_w^{*}$ e o transporte (\code{naturais.tex}).

\textbf{OPERAÇÃO.} Soma com transporte:
\[
a+b \;=\; (a\oplus b) \;+\; 2(a\wedge b).
\]

\textbf{DEMONSTRAÇÃO.} Medido em \code{naturais thm:transporte}; a cadeia
de papers não redefine o que recebe.

\textbf{VOLTA.} Dual da torre: $\nu\circ\nu=\mathrm{id}$.

\section{tick inteiros --- o sinal}\label{sec:tick-Z}
\textbf{DEFINIÇÃO.} $\mathbb{Z}=\mathbb{N}^{2}/{\sim}$ --- o sinal é a
escolha de metade (\code{inteiros.tex}).

\textbf{OPERAÇÃO.} $a+x=b$ tem solução em $\mathbb{Z}$ sempre; em $\mathbb{N}$
só se $b\ge a$.

\textbf{VOLTA.} $\nu(z)=-z$.

\section{tick racionais --- a fibra}\label{sec:tick-Q}
\textbf{DEFINIÇÃO.} $\mathbb{Q}$ --- produto com inverso nos não-nulos
(\code{racionais.tex}).

\textbf{TRADUÇÃO.} \verb|\frac{1}{2}+\frac{1}{3}| $\to$ \texttt{(1)/(2)+(1)/(3)}
$\to$ veste $\frac{5}{6}$.

\textbf{OPERAÇÃO.} $\frac{a}{b}+\frac{c}{d}=\frac{ad+bc}{bd}$ (corpo).

\textbf{VOLTA.} Inverso $\frac{b}{a}$ se $a\neq0$.

\section{tick reais --- o corte}\label{sec:tick-R}
\textbf{DEFINIÇÃO.} O real é o \textbf{corte}, nunca um decimal
(\code{reais.tex}).

\textbf{OPERAÇÃO.} Completude: cinco formas equivalentes (problema~25 da casa).

\textbf{VOLTA.} Cauchy$(\mathbb{Q})/{\sim}$ --- o mesmo ponto por quatro portas
(\code{identifica.h}).

%----------------------------------------------------------------------
\section{tick dragao --- Heighway e G}\label{sec:tick-dragao}
\textbf{DEFINIÇÃO.} $\pi\colon I\to\mathbb{Z}^{2}$ (Heighway);
$G(x)=|\pi^{-1}(x)|$. Auto-intersecção $\Rightarrow G>1$
(\code{arquitetura.tex}, \code{tests/aranha\_g.c}).

\textbf{OPERAÇÃO (aranha).}
\begin{enumerate}\itemsep2pt
\item escreve $G\mathrel{+}=1$;
\item lê vizinhança;
\item decide por $\min G$;
\item fecha $R^{4}=\mathrm{id}$.
\end{enumerate}

\textbf{DEMONSTRAÇÃO.} §AG1--§AG11: Heighway dobra; recta $G\le1$; paridade
$G\bmod2$ em $\mathrm{GF}(2)$.

\textbf{VOLTA (inversa).} $\tilde\pi(i)=(\pi(i),k(i))$, $\tilde G\equiv1$,
$\mathrm{pr}_{1}\circ\tilde\pi=\pi$ (\code{thm:aranha-inversa}).

\section{tick orbita de uma frase}\label{sec:tick-frase}
\textbf{DEFINIÇÃO.} Frase do corpus $\to$ álgebra do idioma ($I$ em
$\Word_8$) $\to$ $\pi$ Heighway $\to$ $\tilde\pi$.

\textbf{OPERAÇÃO.} \code{tests/idioma\_orbita.c} §IO1--§IO5.

\textbf{VOLTA.} Indexar: \code{idioma/pt/orbita/\ldots} via
\code{IMPORT IDIOMA} + \code{corpus\_orbitas\_pt.txt}.

%----------------------------------------------------------------------
\section{1 mais 1 em dois corpos}
Depende do corpo (tick~bit vs tick~$\mathbb{Z}$):
\[
1+1 \;=\;
\begin{cases}
0 & \text{em }\mathrm{GF}(2),\\
2 & \text{em }\mathbb{Z}.
\end{cases}
\]
Sem o corpo declarado a pergunta está incompleta.
ASCII: \texttt{1+1} desdobra na fita; LaTeX devolve os dois casos vestidos.

\section{resolve um terco}
Em base~$10$: $0{,}333\ldots$ (não fecha). Em base~$3$: $0{,}1$ exacto.
Em LaTeX: $\frac{1}{3}$. ASCII desdobrado: \texttt{(1)/(3)}.

\section{resolve a raiz de 2}
Em $\mathbb{Q}$: não existe ($p^{2}=2q^{2}$ força paridade). Em
$\mathbb{Z}/7$: $3^{2}\equiv2$. Em $\mathbb{R}$: o corte positivo.
LaTeX: $\sqrt{2}$; ASCII: \texttt{raiz (2)}.

\section{o que e o dragao}
A curva de Heighway: realização geométrica que dobra --- $G>1$ nas
auto-intersecções. Papers: \code{arquitetura.tex} §aranha;
medidor \code{tests/aranha\_g.c}. Inversa: $\tilde G\equiv1$.

\section{ensina a cadeia}
Começa no bit (\S\ref{sec:tick-bit}), sobe naturais--reais, realiza no
dragão (\S\ref{sec:tick-dragao}). Cada degrau tem paper próprio; este
documento só ordena os \textbf{ticks} e a membrana. Pergunta
«mostra a fundação» para a escada Corpos$\to$Peano
(\code{torre\_fundacao.tex}).

\section{mostra a ciencia do dragao}
Este ficheiro: \code{corpus/docs/ciencia\_dragao.tex}. Índice vivo:
\code{bash tools/indexa\_tex\_vivo.sh}. Resposta em LaTeX ao vivo ---
alterar o \texttt{.tex} altera o que a assistente ensina, sem reingerir.

%======================================================================
\section{o que e o tick}\label{sec:tick-teoria}
Três leituras, \textbf{uma} estrutura --- o quantum do relógio da casa.

\subsection*{Maestro (corpo de Peano)}
O Teorema do Maestro (\code{corpo\_topologico thm:proj-maestro}):
\[
\boxed{\;\mathcal{P}_k \;=\; \mathrm{tick}\circ\mathrm{batuta}\circ\Pi\;}
\]
--- projecta $\pi_k$ a cada tick; \textbf{não} inventa estrutura.
Recebe $(\Pi,\,\mathrm{tick},\,I)$ com batuta $I\in\{-1,0,+1\}$.
O tick é o \textbf{ciclo de processamento} (Lei~6 / janela): uma
passagem do projector. O Metrónomo \emph{lê} o resíduo; o Maestro
\emph{avança} o índice.

\subsection*{Quantizador / relógio (assistente)}
Em \code{conversa.c}: o tick é o quantum da fita --- uma ou $2^{m}$
dobras por chamada (\texttt{TICKS}$\in\{1,2,4,8\}$, andares da torre;
o $3$ recusa-se). O estado mora em disco (\texttt{<base>.tick}): fechar
o programa não perde a conta. \emph{Medido:} a velocidade só sobe por
dobra --- o mesmo princípio $\omega_{2M}^{2}=\omega_M$ do Quantizador
(\code{corpo\_algebrico thm:batuta-continuo}: o gap entre dois ticks
consecutivos \emph{é} o comprimento do segmento afim).

\subsection*{Ensino (seis slots)}
Cada resolução nomeia o corte: DEFINIÇÃO, TRADUÇÃO, OPERAÇÃO,
DEMONSTRAÇÃO, CONTRAEXEMPLO, VOLTA --- \texttt{tique7} no conversa.
Aqui o tick não é tempo físico: é \textbf{granularidade da prova}.

\medskip\noindent\textbf{Síntese.}
\[
\boxed{
\mathrm{tick}
\;=\;
\text{quantum do índice}
\;=\;
\text{uma dobra da torre / uma passagem de }\mathcal{P}_k
\;=\;
\text{um passo nomeado da resolução}.
}
\]

%----------------------------------------------------------------------
\section{regua do idioma --- dinamica}\label{sec:regua-idioma}
A «régua» do idioma \textbf{não} é métrica euclidiana: é a álgebra sobre
$\Word_8$ que \emph{emite} a órbita discreta $I$
(\code{lib/portugues.h}, gémeas EN/ES).

\subsection*{Geração de $I$}
\begin{enumerate}\itemsep2pt
\item $\Sigma\subseteq\Word_8$ --- alfabeto (bytes que existem);
\item monoide $\Sigma^{*}$ --- concatena;
\item léxico $\nu$ com $\nu\circ\nu=\mathrm{id}$ --- dual;
\item regras (concat, dual, identifica) --- transformações no monoide.
\end{enumerate}
\code{pt\_emite\_orbita}: cada byte $\in\Sigma$ é um \textbf{índice};
cada dual léxico fechado é outro índice (código $>255$). Donde
\[
|I| \;=\; \#\{\text{bytes lidos}\} \;+\; \#\{\text{regras aplicadas}\}.
\]
$|I|$ parametriza Heighway: $\pi$ tem $|I|$ pontos
(\code{tests/idioma\_orbita.c}).

\subsection*{Dinâmica}
O tempo do idioma é o índice $i\in I$:
\[
i \;\longmapsto\; i+1
\qquad\text{(ler o próximo byte ou aplicar a próxima regra)}.
\]
Não há relógio externo: a frase \emph{é} a fita; o tick da régua é
avançar $i$. A batuta do Maestro, neste regime, escolhe o sentido
(ler / dual / identificar); o Metrónomo confere que o passo ainda está
em $\Sigma$ ou na regra --- senão o símbolo «não existe no sistema».

\subsection*{Ligação à cifra}
A mesma casa compara por fração contínua (\code{sql.c} / cifra): a
régua do idioma emite $I$; a cifra endereça textos no banco. Duas
réguas, um suporte $\Word_8$ --- como ASCII $\leftrightarrow$
$\mathbb{P}^{1}(\mathbb{F}_{127})$ (\code{lib/ascii.h}).

%----------------------------------------------------------------------
\section{algoritmo inverso e o tick}\label{sec:tick-inverso}
A cadeia \emph{directa} (thm:multiplicidade):
\[
I \;\xrightarrow{\ \pi\ }\; \mathbb{Z}^{2},\qquad
G(x)=|\pi^{-1}(x)|
\quad(G>1\text{ nas dobras Heighway}).
\]
A cadeia \emph{inversa} (thm:aranha-inversa):
\[
\tilde\pi(i)=\bigl(\pi(i),\,k(i)\bigr),\qquad
\tilde G\equiv1,\qquad
\mathrm{pr}_{1}\circ\tilde\pi=\pi.
\]
\textbf{Cada $i\in I$ é já um tick} da realização. O inverso não cria
tempo novo: em cada tick $i$,
\begin{enumerate}\itemsep2pt
\item lê $x=\pi(i)$ --- estado na grade;
\item incrementa $c[x]$; põe $k(i)=c[x]$ --- folha;
\item emite $\tilde\pi(i)=(x,k(i))$ --- multiplicidade~$1$.
\end{enumerate}
A fibra que $|\det|=1$ perdia (§M10) restaura-se \emph{tick a tick}
como coordenada $k$.

%----------------------------------------------------------------------
\section{integrar o tick no algoritmo}\label{sec:tick-integra}
Uma só linha de tempo para idioma + dragão + Maestro:

\[
\boxed{\;
\begin{array}{c}
\text{frase}
\;\xrightarrow{\ \mathrm{tick}_{alg}\ }\;
I[i]
\;\xrightarrow{\ \mathrm{tick}_{\pi}\ }\;
\pi(i)
\;\xrightarrow{\ \mathrm{tick}_{\tilde\pi}\ }\;
\tilde\pi(i)
\\[0.6em]
\mathcal{P}
\;=\;
\mathrm{tick}\circ\mathrm{batuta}\circ\Pi
\quad\text{com}\quad
\Pi=\text{frase/partitura},\;
I_{\mathrm{bat}}\in\{-1,0,+1\}.
\end{array}
\;}
\]

\begin{center}\small
\begin{tabular}{@{}lll@{}}
\toprule
papel & no idioma/dragão & no Maestro \\
\midrule
$\Pi$ & frase / partitura & especificação \\
batuta $I$ & ler / dual / identificar & trial $\{-1,0,+1\}$ \\
tick & avançar $i\in I$ & ciclo $\mathcal{P}_k$ \\
$\pi_k$ / $\pi$ & Heighway no passo $i$ & retração de andar \\
Metrónomo & $\tilde G\equiv1$, resíduo $0$ & auditoria / Lyapunov \\
\bottomrule
\end{tabular}
\end{center}

\noindent\textbf{Consequência operacional.}
\begin{enumerate}\itemsep2pt
\item \texttt{pt\_emite\_orbita} --- produz a fita de ticks algébricos;
\item Heighway --- um tick geométrico por índice (escreve $G\mathrel{+}=1$);
\item inverso --- um tick de levantamento por índice;
\item assistente --- \texttt{tick "\ldots"} / \texttt{tique7} nomeiam o
      mesmo quantum na fala (Quantizador + ensino).
\end{enumerate}
Não são quatro relógios: é \textbf{um} índice $i$, lido em quatro
cartas (álgebra, grade, folha, discurso). Integrar o tick no algoritmo
é fazer o Maestro avançar $i$ sobre a órbita do idioma, e o Metrónomo
exigir $\tilde G\equiv1$ (ou resíduo~$0$ na fita) a cada passo.

\medido{\code{tests/tick\_idioma.c} §TI1--§TI6 ($6{:}0$): $\Pi$ emite $I$
via \code{pt\_emite\_orbita}; o Maestro avança um tick $\to$ um ponto
Heighway; batuta $+1/-1$ partilha $G$ e a ordem é dual; levantamento
tick~a~tick com $\tilde G\equiv1$; $P=\mathrm{tick}\circ\mathrm{batuta}\circ\Pi$
coincide com Heighway em lote; Metrónomo $\mathrm{pr}_{1}\circ\tilde\pi=\pi$,
$\sum G=|I|$.}

\section{o tick na teoria do dragao}
Resposta curta: o tick é o quantum do índice $i\in I$; a régua do
idioma emite $I$; Heighway realiza; o inverso levanta a fibra
tick~a~tick; o Maestro é $\mathrm{tick}\circ\mathrm{batuta}\circ\Pi$
sobre essa mesma fita. Fontes:
\code{corpo\_topologico thm:proj-maestro},
\code{arquitetura thm:aranha-inversa},
\code{lib/portugues.h},
\code{conversa.c} (relógio / \texttt{TICKS}).

\end{document}
